\section{Half-integer sums and their differential system}\label{sec:forms}

We now turn the rational function of Proposition~\ref{prop:rational}
into linear forms in beta values. First we identify the functions whose
values at $-1$ give the required even arguments. We then express the sums
of $F_n$ in terms of those functions and remove the finite corrections
by differentiation.

\subsection{The functions at the evaluation point}

The half-integer denominators are chosen so that alternating signs give
the beta series: $(t+\tfrac12)^{-j}=2^j(2t+1)^{-j}$.
For $j\geq1$, put
\[
 \Phi_j(z)=\sum_{t=0}^{\infty}\frac{z^t}{(t+\tfrac12)^j}
 \qquad(|z|<1).
\]
The integral
\begin{equation}\label{eq:Phi-integral}
 \Phi_j(z)=\frac1{(j-1)!}\int_0^1
       \frac{x^{-1/2}(-\log x)^{j-1}}{1-zx}\dd x
\end{equation}
continues it to $\C\setminus[1,\infty)$.
In particular, the series or its Abel limit gives
$\Phi_j(-1)=2^j\beta(j)$, including $j=1$.

We still need to remove the odd arguments $j$. The positive and negative
half-integer sums lead to functions at $z$ and $1/z$; we combine them as
\[
 H_j(z)=z^{-1}\Phi_j(1/z)-(-1)^j\Phi_j(z).
\]
Both terms are analytic near $-1$ by \eqref{eq:Phi-integral}.
At that point the combination supplies the factor $1+(-1)^j$:
\begin{equation}\label{eq:H-values}
 H_j(-1)=-(1+(-1)^j)2^j\beta(j).
\end{equation}
Thus $H_j(-1)$ is zero for odd $j$ and is $-2^{j+1}\beta(j)$ for even $j$.

Their derivatives are rational-function combinations of the same
functions and the constant function. Let $\Phi_0(z)=1/(1-z)$.
For any positive integer $d$, the identity
$z\Phi_j'(z)=\Phi_{j-1}(z)-\Phi_j(z)/2$ implies that
$(1,H_1,\ldots,H_d)$ is a solution of
\begin{align}
 y_0'&=0,\notag\\
 y_1'&=\frac{2}{z(1-z)}y_0-\frac1{2z}y_1,\label{eq:base-system}\\
 y_j'&=-\frac1z y_{j-1}-\frac1{2z}y_j \qquad(2\leq j\leq d).\notag
\end{align}
The coordinate $y_0$ represents the constant function $1$, which appears
in the derivative of $H_1$. Since the coefficients have no pole at $-1$,
arbitrary initial values there determine an analytic solution. We will
use the beta-valued initial vector for the linear forms, and arbitrary
initial vectors for the rank proof.

\subsection{Summation and finite tails}

From now on we assume $r\geq2$, and set $M=rn$.
The sums start beyond this cutoff so that the decay estimates for $F_n$
apply. To obtain several families of forms we use $F_n$ and its first
$h$ derivatives, where $0\leq h\leq a$ is an integer. A derivative of
order $p$ raises a denominator power $i$ to $i+p$, so we use the system
\eqref{eq:base-system} with $d=a+h$.

The powers of $z$ in the following definition align the two reflected
sums with $H_j$. The factor $z^M$ makes the finite corrections polynomial,
as the next lemma shows.
For $0\leq p\leq h$ define, on the unit circle,
\begin{equation}\label{eq:S}
 S_{n,p}(z)=z^M\sum_{t=M}^{\infty}
 \left\{F_n^{(p)}(t+\tfrac12)z^{-t-1}
              -F_n^{(p)}(-t-\tfrac12)z^t\right\}.
\end{equation}
Here $F_n^{(p)}$ is differentiated with respect to its argument.
Near $-1$ we continue the two tails separately, using the identity below.

\begin{lemma}\label{lem:tail}
Near $-1$ there is an identity
\begin{equation}\label{eq:tail-decomposition}
 S_{n,p}(z)=V_{n,p}(z)
      +z^M\sum_{i=1}^b(-1)^p(i)_p P_i(z)H_{i+p}(z),
 \qquad \deg V_{n,p}\leq2M-1.
\end{equation}
\end{lemma}

\begin{proof}
Differentiate the partial fractions of $F_n$. For a fixed pair $i,j$
write $s=i+p$ and change the summation index to $l=t+j$ in the positive
tail, and to $l=t-j$ in the negative tail. The infinite parts give the
corresponding term of \eqref{eq:tail-decomposition}. The finite correction
is $(-1)^p(i)_p c_{i,j}$ times
\begin{equation}\label{eq:finite-correction}
 -\sum_{l=0}^{M+j-1}\frac{z^{M+j-l-1}}{(l+\tfrac12)^s}
 +(-1)^s\sum_{l=0}^{M-j-1}\frac{z^{M+j+l}}{(l+\tfrac12)^s}.
\end{equation}
Since $0\leq j\leq n\leq M$, all exponents are nonnegative and at most
$2M-1$. Summing these corrections proves the assertion.

The positive and negative infinite tails are initially analytic on
opposite sides of the circle. Each has the displayed continuation near
$-1$ by \eqref{eq:Phi-integral}. For the derivative orders used below,
the differentiated series converge absolutely on the circle; their
values agree with the derivatives of the continuations by Abel limits.
\end{proof}

\subsection{Rows representing derivatives}

Write the nonpolynomial part of \eqref{eq:tail-decomposition} as
$\sum_i C_i^{[p]}H_i$. Its coefficient polynomials are
\begin{equation}\label{eq:C}
 C_0^{[p]}=0,\qquad
 C_i^{[p]}(z)=
 \begin{cases}
 (-1)^p(i-p)_p z^M P_{i-p}(z),&p+1\leq i\leq p+b,\\
 0,&\text{otherwise}.
 \end{cases}
\end{equation}
The degree of each is at most $M+n$.
They have integer coefficients. We next track the coefficients produced
by differentiating this expression. The system \eqref{eq:base-system}
expresses every derivative in the same functions, so no additional
special functions are needed.
Starting with $Q_{1,i}^{[p]}=C_i^{[p]}$, the product rule gives
\begin{align}
 Q_{k+1,0}^{[p]}&=(Q_{k,0}^{[p]})'
                   +\frac{2Q_{k,1}^{[p]}}{z(1-z)},\notag\\
 Q_{k+1,i}^{[p]}&=(Q_{k,i}^{[p]})'
                   -\frac{Q_{k,i}^{[p]}}{2z}
                   -\frac{Q_{k,i+1}^{[p]}}z
                    \qquad(1\leq i\leq d),                    \label{eq:Q}
\end{align}
where $Q_{k,d+1}^{[p]}=0$. The index $k=1$ corresponds to no
differentiation; hence $Q_k^{[p]}$ represents the derivative of order
$k-1$.
For \emph{every} solution of \eqref{eq:base-system}, differentiation gives
\begin{equation}\label{eq:universal-row}
 \left(\sum_{i=1}^d C_i^{[p]}y_i\right)^{(k-1)}
       =\sum_{i=0}^d Q_{k,i}^{[p]}y_i.
\end{equation}
For $k\geq2M+2$, the derivative also kills $V_{n,p}$, so
\begin{equation}\label{eq:value-row}
 \frac{S_{n,p}^{(k-1)}(-1)}{(k-1)!}
  =\frac{Q_{k,0}^{[p]}(-1)}{(k-1)!}
     +\sum_{i=1}^d\frac{Q_{k,i}^{[p]}(-1)}{(k-1)!}H_i(-1).
\end{equation}
Equation \eqref{eq:H-values} makes the right side a rational linear form
in $1$ and the even beta values. Dividing by $(k-1)!$ turns the derivatives
into Taylor coefficients, whose denominators we can control in the next
section. The finite correction $V_{n,p}$ has disappeared, so its
coefficient denominators do not have to be cleared.
The universal identity \eqref{eq:universal-row}, valid for arbitrary
initial values of the solution, will supply the rank argument.

